\documentclass[12pt,a4paper]{report}

% ------------------------------------------------
% PACKAGES
% ------------------------------------------------
\usepackage[margin=1in]{geometry}
\usepackage{graphicx}
\usepackage{float}
\usepackage{booktabs}
\usepackage{array}
\usepackage{longtable}
\usepackage{xcolor}
\usepackage{hyperref}
\usepackage{setspace}
\usepackage{titlesec}
\usepackage{enumitem}

% ------------------------------------------------
% FORMATTING
% ------------------------------------------------
\onehalfspacing

\titleformat{\chapter}
{\normalfont\huge\bfseries}
{\thechapter.}{15pt}{}

\titleformat{\section}
{\normalfont\Large\bfseries}
{\thesection}{10pt}{}

\hypersetup{
    colorlinks=true,
    linkcolor=black,
    urlcolor=blue
}

\begin{document}

% =================================================
% TITLE PAGE
% =================================================

\begin{titlepage}
\thispagestyle{empty}

\begin{center}


{\Large \textbf{SmartExpense: Frontend Prototype Report}\par}

\vspace{0.2cm}

{\normalsize Software Engineering Project\par}

\vspace{0.15cm}

{\normalsize Academic Year 2026--2027\par}

\vspace{0.5cm}

{\Large \textbf{THAPAR INSTITUTE OF ENGINEERING AND}\\
\textbf{TECHNOLOGY}\par}

\vspace{0.4cm}

{\Large \textbf{SmartExpense: Smart Expense Tracker \&}\\
\textbf{Event-Based Savings Planner}\par}

\vspace{0.25cm}

{\normalsize
A Browser-Based Personal Finance Dashboard with\\
Event Savings Goals and Linked Checklists
\par}

\vspace{0.7cm}

\textbf{Submitted by:}

\vspace{0.15cm}

 \quad Divanshi Garg\ [1024030441]
 \vspace{0.15cm}
 \quad Navya Arora [1024030438]

\vspace{0.35cm}

B.E. Third Year, Computer Science \& Engineering

\vspace{0.5cm}

\textbf{Submitted to:}

\vspace{0.15cm}

Dr. Jeelani Asif\\
Department of Computer Science and Engineering

\vspace{1.2cm}

{\normalsize Frontend Prototype Evaluation Report\par}

\vspace{0.1cm}

{\normalsize Computer Science and Engineering Department\par}

\end{center}

\end{titlepage}

\clearpage
% =================================================
% TABLE OF CONTENTS
% =================================================

\tableofcontents



% =================================================
% CHAPTER 1
% =================================================

\chapter{Introduction}

\section{Background and Context}

Keeping track of everyday spending is a problem most people solve badly:
expenses are scattered across cash, UPI, cards and bill payments. By the
time a monthly statement arrives, it may be too late to change spending
behaviour.

This becomes more difficult around specific occasions such as a festival,
wedding, birthday or trip, where a person is not only managing everyday
expenses but is also trying to save towards a fixed and dated goal.
Most existing expense trackers either provide a generic monthly ledger
without linking it to upcoming events or focus only on expense
categorisation.

SmartExpense is proposed as a personal finance web application that
addresses this gap. It provides a single dashboard for monthly spending,
transactions, savings events with target budgets and progress indicators,
and checklists linked to individual events.

\subsection{Problem Statement}

Individuals lack a single lightweight tool that connects ordinary monthly
expense tracking with saving towards specific upcoming events. Spreadsheets
and generic budgeting applications often treat all spending as one stream.

As a result, users may find it difficult to answer questions such as:

\begin{itemize}
    \item How much have I saved towards Diwali?
    \item How much is remaining for an upcoming event?
    \item What items are still left to arrange for a wedding?
\end{itemize}

SmartExpense addresses this problem by combining category-wise expense
tracking, event-based savings goals and linked checklists within one
interface.

\section{Project Objectives}

The core objectives of the SmartExpense project are:

\begin{enumerate}
    \item \textbf{Centralised Spending Dashboard:} Provide a single monthly
    overview of total spending, remaining budget, upcoming events and
    category-wise spending.

    \item \textbf{Transaction Logging:} Allow transactions to be added
    manually or imported through a CSV file.

    \item \textbf{Event-Based Savings Tracking:} Allow users to create
    events with a target date and budget and track savings progress using
    visual progress indicators.

    \item \textbf{Linked Checklists:} Provide event-specific notes and
    checklists so planning tasks remain connected to their related events.

    \item \textbf{Consistent Interface:} Provide a shared sidebar layout
    across Dashboard, Transactions, Events, Notes and Profile screens.

    \item \textbf{Foundation for Full-Stack Development:} Build and validate
    a complete frontend prototype that can later be connected to a backend,
    database and predictive spending service.
\end{enumerate}


% =================================================
% CHAPTER 2
% =================================================

\chapter{Methodology and Design}

\section{System Architecture}

\subsection{High-Level Design}

At the current frontend-only stage, SmartExpense is organised into three
main layers:

\begin{enumerate}
    \item Presentation Layer
    \item Client Logic Layer
    \item Local Data Store
\end{enumerate}

The Presentation Layer consists of static HTML and CSS pages. The Client
Logic Layer is implemented using \texttt{script.js}, which reads and writes
application data. The Local Data Store uses the browser's
\texttt{localStorage}.

The backend Application Layer is planned for a future phase.

% Add architecture image here
\begin{figure}[h]
    \centering
    \includegraphics[width=0.7\textwidth]{ChatGPT Image Sep 14, 2026, 09_32_57 AM.png}
    \caption{My Image}
    \label{fig:myimage}
\end{figure}

The current application flow is as follows:

\begin{enumerate}
    \item A page is loaded.
    \item \texttt{script.js} checks and seeds sample data on the first run.
    \item Relevant helper functions read data from \texttt{localStorage}.
    \item The page render function creates cards, tables or grids.
    \item Add or edit operations update the stored data.
    \item The page is immediately re-rendered.
\end{enumerate}

No network request is made in the current frontend prototype.

\subsection{Technical Stack}

The SmartExpense prototype uses the following technologies:

\begin{enumerate}
    \item \textbf{Markup and Styling:} HTML5 and CSS3.
    \item \textbf{Typography:} Google Fonts, including Playfair Display and Inter.
    \item \textbf{Client Scripting:} Vanilla JavaScript using a shared
    \texttt{script.js} file.
    \item \textbf{Client-Side Persistence:} Browser \texttt{localStorage}.
    \item \textbf{Icons:} Font Awesome and Unicode glyphs.
    \item \textbf{Development Environment:} Standard HTML, CSS and JavaScript
    development tested directly in the browser.
    \item \textbf{Version Control:} Git.
\end{enumerate}


\section{Implementation Details}

\subsection{Module 1: Core Screens}

The core module implements the shared application shell and Dashboard.

\begin{enumerate}
    \item \textbf{Login Screen:} Contains an email and password form with
    client-side required-field validation. Successful submission redirects
    the user to the Dashboard.

    \item \textbf{Shared Sidebar:} Provides navigation links for Dashboard,
    Transactions, Events, Notes and Profile.

    \item \textbf{Dashboard:} Displays summary cards, category-wise spending,
    recent transactions and upcoming events.

    \item \textbf{seedDefaultsIfEmpty():} Populates \texttt{localStorage}
    with realistic sample data on the first run.
\end{enumerate}

\begin{figure}[h]
    \centering
    \includegraphics[width=0.7\textwidth]{pipeline.png}
    \caption{Core page-load and Render pipeline}
    \label{fig:myimage}
\end{figure}


\subsection{Module 2: Secondary Features}

Three additional modules extend SmartExpense to its full frontend scope.

\begin{enumerate}
    \item \textbf{Transactions:} Displays transactions in a sortable table
    with merchant, category, amount, date and entry type. It supports manual
    transaction addition and CSV import.

    \item \textbf{Events:} Displays savings events with target dates,
    days-remaining status, budgets and savings progress bars.

    \item \textbf{Notes:} Provides event-linked checklists with checkable
    items whose completion state is saved in \texttt{localStorage}.

    \item \textbf{Profile:} Allows users to update their name, email and
    monthly budget.
\end{enumerate}

\begin{figure}[H]
    \centering
     \includegraphics[width=0.9\textwidth]{51ec99dd-f4f7-4283-9482-3b4490b1ce16.png}
    
    \caption{Secondary Feature Flow}
    \label{fig:secondaryflow}
\end{figure}


% =================================================
% CHAPTER 3
% =================================================

\chapter{Results and Evaluation}

\section{Testing Methodology}

Since the prototype does not currently contain a backend, testing was
performed entirely in the browser. Each of the six screens was loaded and
tested manually:

\begin{itemize}
    \item Login
    \item Dashboard
    \item Transactions
    \item Events
    \item Notes
    \item Profile
\end{itemize}

The browser DevTools console and DOM inspector were used to check that
elements populated correctly and that no JavaScript errors occurred.
The browser's \texttt{localStorage} was also inspected after add and edit
operations to confirm that application data was stored correctly.

\begin{figure}[H]
    \centering
     \includegraphics[width=0.85\textwidth]{3e3b5047-20ad-4636-b620-88e63f1c3c3a.png}
   
    \caption{Manual Page-by-Page Testing Framework}
    \label{fig:testing}
\end{figure}


\section{Feature Completion Status}

Table \ref{tab:features} summarises the current status of the planned
SmartExpense features.

\begin{table}[H]
\centering
\caption{SmartExpense Feature Completion Status}
\label{tab:features}

\renewcommand{\arraystretch}{1.4}

\begin{tabular}{|p{5cm}|p{5cm}|p{3cm}|}
\hline
\textbf{Capability} & \textbf{Current State} & \textbf{Status} \\
\hline

Dashboard summary cards and category chart &
Working from localStorage &
Done (Frontend) \\
\hline

Add transaction / CSV import &
Working with prompt-based UI &
Done (Frontend) \\
\hline

Event creation and savings progress &
Working with prompt-based UI &
Done (Frontend) \\
\hline

Event-linked checklists &
Working, toggle persists &
Done (Frontend) \\
\hline

Profile and budget editing &
Working, writes to profile &
Done (Frontend) \\
\hline

User authentication / sessions &
Form only, no backend verification &
Pending \\
\hline

Persistent multi-device storage &
Single-browser localStorage only &
Pending \\
\hline

Predicted spending forecast &
Placeholder card in UI &
Pending \\
\hline

Spending anomaly alert &
Markup present but not connected &
Pending \\
\hline

\end{tabular}
\end{table}


\section{Analysis}

At the current stage, the frontend prototype demonstrates the intended
SmartExpense user experience, including dashboard overview, transaction
logging, event-based savings and linked checklists.

The prototype confirms that the main data model consisting of transactions,
events, notes and profile information can support the application's main
features.

The primary remaining work includes:

\begin{itemize}
    \item Replacing \texttt{localStorage} with a backend and database.
    \item Adding real user authentication.
    \item Implementing persistent multi-device storage.
    \item Connecting the predicted-spending feature to an analysis service.
    \item Implementing spending anomaly detection.
\end{itemize}


% =================================================
% CHAPTER 4
% =================================================

\chapter{Conclusion}

\section{Summary of Work}

This project focused on building SmartExpense, a web application that allows
users to track everyday spending and save towards specific upcoming events
from a single dashboard.

A frontend-only prototype was designed and implemented with the following
screens:

\begin{itemize}
    \item Login
    \item Dashboard
    \item Transactions
    \item Events
    \item Notes
    \item Profile
\end{itemize}

These screens share a client-side data model stored in browser
\texttt{localStorage}. The prototype was validated using realistic sample
data and manual testing of each screen and interaction.

\section{Key Findings}

The major findings of the project are:

\begin{enumerate}
    \item A client-side model containing transactions, events, notes and
    profile information is sufficient to represent the complete frontend
    SmartExpense experience.

    \item Using one shared \texttt{script.js} file allows helper functions
    for formatting, date calculations and localStorage operations to be
    reused across multiple pages.

    \item Using \texttt{localStorage} during the prototype stage was an
    effective way to validate the complete user interface before implementing
    a backend database.
\end{enumerate}


\section{Future Work and Recommendations}

The next phase of the SmartExpense project will include:

\begin{enumerate}
    \item Development of a backend REST API.
    \item Integration of a persistent database.
    \item Real user authentication with sign-up, login and logout.
    \item Multi-device data access.
    \item A predictive spending service.
    \item Spending anomaly detection.
    \item Server-side validation for transactions and CSV imports.
    \item Replacing browser \texttt{prompt()} dialogs with proper forms
    and modals.
    \item Load, response-time and usability testing.
\end{enumerate}


\section{Final Remarks}

SmartExpense addresses the problem of everyday spending and event-based
saving being managed through separate and disconnected tools. The frontend
prototype demonstrates that these features can be integrated into one
coherent interface.

The completed prototype provides a dashboard, transaction ledger,
event-based savings tracking and linked checklists using a shared
client-side data model. The next major step is to connect this validated
interface to a real backend and evaluate the complete system at a realistic
scale.

% =================================================
% END DOCUMENT
% =================================================

\end{document}
